\section{Related Work}

\subsection{Theme 1: Deep Learning and Neural Networks for Theorem Proving}

Several studies explore using deep learning and neural networks for theorem proving, developing specialized environments like GamePad and HOList \cite{huang2019gamepad} \cite{bansal2019holist}, and applying techniques like generative language modeling and neural automated theorem proving \cite{polu2020generative} \cite{jiang2023draft}. However, these approaches assume fitting deep learning into traditional theorem proving frameworks \cite{li2021isarstep}. This research proposes enhancing large language models (LLMs) directly with formal mathematical proofs from interactive theorem provers, representing a shift toward independently advancing LLMs' formal reasoning abilities.

\subsection{Theme 2: Formalizing and Mechanizing Mathematical Proofs}

Converting informal proofs to formal versions is explored, with neural networks used for understanding mathematical context \cite{wu2022autoformalization}, leveraging datasets like Lean's proof library \cite{yang2023leandojo}, and utilizing meta-languages \cite{jiang2022thor} and elaboration models from proof assistants \cite{szegedy2020promising}. However, LLMs can reduce bottlenecks by utilizing their knowledge of informal proofs \cite{azerbayev2023llemma} and training on formal proofs to improve automated reasoning \cite{zheng2022minif2f}.

\subsection{Theme 3: Advanced Machine Learning for Theorem Proving}

Various methodologies integrate language models \cite{han2022proof}, self-supervised training \cite{wu2022autoformalization_arxiv}, learning inductive biases \cite{zheng2022minif2f}, curriculum learning \cite{polu2022formal}, and benchmarks like MATH \cite{hendrycks2021measuring} for enhancing mathematical reasoning and theorem proving. In contrast, this research directly leverages verified proofs from interactive theorem provers like Lean to enhance LLMs' reasoning, autoformalization abilities, and evaluation via comprehensive metrics like MATH, ProofNet, and MMLU.
